% ================================================================
% Chapter 14: Ideatic Auctions and Ritual Economy
% ================================================================

\chapter{Ideatic Auctions and Ritual Economy}
\label{ch:auction}

\epigraph{Conspicuous consumption of valuable goods is a means of reputability
to the gentleman of leisure.}{Thorstein Veblen,
\textit{The Theory of the Leisure Class} (1899)}

% ----------------------------------------------------------------
\section{Introduction}
% ----------------------------------------------------------------

In many cultural contexts, ideas compete for interpretive priority through
\emph{depth-based bidding}.  The potlatch of the Pacific Northwest Coast, in
which rival chiefs destroy or give away enormous quantities of wealth to
establish precedence, is the paradigmatic example: the deepest cultural
contribution wins prestige.  But the pattern is far more widespread.
The competitive almsgiving of medieval Christendom, the monumental
architecture of Egyptian pharaohs, the conspicuous sacrifice in ancient
temple economies, the academic citation game---all are forms of
``auction'' where the most complex (deepest) offering wins.

Veblen's theory of conspicuous consumption captures the demand side of this
dynamic: prestige accrues to the display of costly (deep) goods.
Bourdieu's concept of symbolic capital captures the supply side: cultural
producers bid for status by offering complex cultural products to
audiences whose capacity to appreciate them (their cultural capital)
constrains the impact of even the most extravagant bid.\footnote{Bourdieu,
\textit{La Distinction} (1979) and \textit{Les Règles de l'art} (1992).
The ``field of cultural production'' is precisely our auction, with
depth as the currency and resonance as the condition of successful
reception.}

This chapter formalizes ideatic auctions.  We introduce \emph{bid value}
(depth), a \emph{maxDepthIdea} selector, and \emph{auction interpretation}
(the winning bid composed with the receiver's background).  We prove
eighteen theorems showing that void always loses, deeper bids win,
composing bids is subadditive, and the winner's interpretation is bounded
by audience capacity.

% ----------------------------------------------------------------
\section{Core Definitions}
\label{sec:ch14-defs}
% ----------------------------------------------------------------

\begin{definition}[Bid Value]
\label{def:bid-value}
The \emph{bid value} of an idea $a$ is its depth:
$\operatorname{bidValue}(a) = \depth(a)$.
In the ritual economy, depth equals prestige equals bid amount.
\end{definition}

\begin{definition}[Maximum Depth Idea]
\label{def:max-depth}
Given a fallback idea and a list of bids, $\operatorname{maxDepthIdea}$
selects the idea with the greatest depth via a left fold that keeps the
deeper of each pair:
\[
  \operatorname{maxDepthIdea}(\text{best}, []) = \text{best}, \qquad
  \operatorname{maxDepthIdea}(\text{best}, x :: \text{rest}) =
  \begin{cases}
    \operatorname{maxDepthIdea}(x, \text{rest}) & \text{if } \depth(x) \ge \depth(\text{best}), \\
    \operatorname{maxDepthIdea}(\text{best}, \text{rest}) & \text{otherwise}.
  \end{cases}
\]
\end{definition}

\begin{lstlisting}[caption={Maximum depth idea selector in Lean~4}]
def maxDepthIdea : I -> List I -> I
  | best, []       => best
  | best, x :: rest =>
    if depth x >= depth best then maxDepthIdea x rest
    else maxDepthIdea best rest
\end{lstlisting}

\begin{definition}[Auction Interpretation]
\label{def:auction-interp}
The \emph{auction interpretation} composes the receiver's background with
the winning (deepest) bid:
\[
  \operatorname{auctionInterpret}(\text{bids}, \text{fallback}, \text{receiver})
  = \text{receiver} \compose \operatorname{maxDepthIdea}(\text{fallback}, \text{bids}).
\]
The receiver's background mediates the winning bid: the audience does not
absorb the winning idea raw but interprets it through their own cultural
lens.
\end{definition}

% ----------------------------------------------------------------
\section{Void as Non-Competitor}
\label{sec:ch14-void}
% ----------------------------------------------------------------

\begin{theorem}[Void Bid Value is Zero (14.1)]
\label{thm:void-bid}
The void idea has zero bid value: $\operatorname{bidValue}(\void) = 0$.
\end{theorem}

\begin{proof}[Proof sketch]
$\depth(\void) = 0$ by the void depth axiom.
\end{proof}

In the potlatch, offering nothing is the ultimate humiliation.  Zero depth
equals zero cultural capital equals zero bidding power.  The formal basis of
``losing face'': the chief who brings no gifts has bid $\void$ and receives
zero prestige.  Mauss's \emph{obligation to give} is, at bottom, an
obligation to bid non-void.\footnote{Mauss, \textit{Essai sur le don}
(1925): ``To refuse to give, to fail to invite, just as to refuse to
accept, is tantamount to declaring war; it is to reject the bond of
alliance and commonality.''}

\begin{theorem}[Void Loses to Any Bid (14.2)]
\label{thm:void-loses}
Given any bid $a$ and void as fallback, the maximum has non-negative depth:
$\operatorname{bidValue}(\operatorname{maxDepthIdea}(\void, [a])) \ge 0$.
\end{theorem}

\begin{proof}[Proof sketch]
Depth is always non-negative.
\end{proof}

Any cultural offering, no matter how shallow, beats silence.  This is the
formal foundation of ``showing up'': participation is always at least as
good as absence in competitive ritual.  The anthropological wisdom that
``the worst thing you can do is not come'' is a theorem.

\begin{theorem}[Deeper Beats Void (14.3)]
\label{thm:deeper-void}
If $\depth(a) \ge 1$, then $a$ beats void:
$\depth(\operatorname{maxDepthIdea}(\void, [a])) \ge 1$.
\end{theorem}

\begin{proof}[Proof sketch]
When $\depth(a) \ge 1 \ge 0 = \depth(\void)$, the selector chooses $a$.
\end{proof}

In competitive gift-giving, the party who brings \emph{something} always
outranks the party who brings nothing.  Mauss's obligation to give is
structurally enforced: non-givers are structurally dominated.  Even the
most modest gift---a single shell necklace, a handshake, a word of
greeting---outbids silence.

% ----------------------------------------------------------------
\section{Selection Properties}
\label{sec:ch14-selection}
% ----------------------------------------------------------------

\begin{theorem}[Single Bidder Wins (14.4)]
\label{thm:single-bidder}
With only one bid, the winner is determined by comparison with the fallback:
\[
  \operatorname{maxDepthIdea}(\text{fallback}, [a]) =
  \begin{cases}
    a & \text{if } \depth(a) \ge \depth(\text{fallback}), \\
    \text{fallback} & \text{otherwise}.
  \end{cases}
\]
\end{theorem}

\begin{proof}[Proof sketch]
Direct unfolding of the definition with an empty remainder.
\end{proof}

A monopoly on cultural production---single priesthood, state media---means
the sole producer always ``wins'' the auction (assuming their depth exceeds
the fallback tradition).  Competition is absent; interpretation is
determined unilaterally.  This is the formal structure of theocracy:
one bidder, guaranteed victory.

\begin{theorem}[Fallback Survives Empty Auction (14.5)]
\label{thm:empty-auction}
With no bids, the fallback wins:
$\operatorname{maxDepthIdea}(\text{fallback}, []) = \text{fallback}$.
\end{theorem}

\begin{proof}[Proof sketch]
The base case of the recursion.
\end{proof}

When no one competes, the status quo persists.  The ``default culture''
(tradition, orthodoxy) wins by forfeit when no challenger appears.  This
is the formal mechanism of cultural conservatism: in the absence of
competing bids, the existing cultural framework is preserved intact.

\begin{theorem}[Bid Value is Non-Negative (14.6)]
\label{thm:bid-nonneg}
All bids have non-negative value: $\operatorname{bidValue}(a) \ge 0$.
\end{theorem}

\begin{proof}[Proof sketch]
Depth is a natural number.
\end{proof}

There are no ``negative contributions'' in ideatic space.  Every idea has
non-negative complexity.  Even the most trivial cultural offering---a
monosyllable, a blank canvas---has at least zero depth.  The floor of
cultural value is zero, never negative.  This is a modeling choice with
significant anthropological implications: we assume that ideas cannot
\emph{destroy} depth, only fail to add it.\footnote{One could extend the
model with signed depth to capture ``anti-ideas'' or cultural
destruction, but the current axiomatization does not permit this.  The
non-negativity of depth is one of IDT's stronger structural commitments.}

% ----------------------------------------------------------------
\section{Composition and Bidding}
\label{sec:ch14-composition}
% ----------------------------------------------------------------

\begin{theorem}[Composed Bid Subadditivity (14.7)]
\label{thm:composed-bid}
The bid value of composing two ideas is at most the sum of their bid values:
\[
  \operatorname{bidValue}(a \compose b)
  \;\le\;
  \operatorname{bidValue}(a) + \operatorname{bidValue}(b).
\]
\end{theorem}

\begin{proof}[Proof sketch]
By subadditivity of depth.
\end{proof}

Combining two gifts into one package does not create more prestige than
giving them separately.  The potlatch chief who presents a composite gift
of blankets and copper shields gains at most the sum of their separate
values.  Subadditivity prevents ``prestige inflation'' through bundling.
This is the formal refutation of the marketing claim that ``the whole is
greater than the sum of its parts''---in ideatic auctions, the whole is
at most equal to the sum.

\begin{example}[The Contemporary Art Market]
\label{ex:art-market}
The contemporary art market, epitomized by the auction houses Sotheby's
and Christie's, provides a nearly literal instantiation of the ideatic
auction.  When a painting by, say, Jean-Michel Basquiat is offered at
auction, each bidder places a monetary bid that functions as a proxy for
their valuation of the work's cultural depth.  The winning bidder---the
one who values the work most highly (or who has the most resources to
signal status)---acquires the right to ``interpret'' the work: to display
it, to incorporate it into their collection narrative, to claim the
prestige associated with ownership.

Theorem~\ref{thm:composed-bid} applies directly: the prestige of
displaying two works together (a Basquiat alongside a Warhol) is at most
the sum of their individual prestige values.  Bundling artworks does not
create prestige beyond the additive bound.  Theorem~\ref{thm:void-bid}
explains why the artist who produces nothing (the rumored but never
delivered masterpiece) commands zero prestige in the market, no matter
how great their reputation.  And Theorem~\ref{thm:winner-bound} captures
a well-known phenomenon in art criticism: the depth of an artwork's
reception depends as much on the audience's sophistication ($\depth(r)$)
as on the work itself.  A Rothko color field painting that elicits
profound meditative experience in a viewer steeped in Abstract
Expressionism may produce only bewilderment in a viewer without that
background.\footnote{For the sociology of the art market, see Don Thompson,
\textit{The \$12 Million Stuffed Shark} (2008); and Olav Velthuis,
\textit{Talking Prices: Symbolic Meanings of Prices on the Market for
Contemporary Art} (2005).}
\end{example}

\begin{remark}[Connection to Mechanism Design]
\label{rem:mechanism-design}
The ideatic auction framework connects to the broader theory of
\emph{mechanism design}---the ``reverse engineering'' of game theory, in
which the designer chooses the rules of the game to achieve a desired
outcome.  In our setting, the ``mechanism'' is the cultural institution
that determines how bids are solicited, evaluated, and rewarded.  The
potlatch is one mechanism; the academic peer review system is another;
the art auction is a third.  Each implements a different set of rules
for depth-based competition, but all are subject to the same structural
constraints: void bids lose, composition is subadditive, and reception
is bounded.

The mechanism design perspective suggests a normative question: what
cultural institutions \emph{should} we design?  If the goal is to maximize
the depth of the winning bid (cultural excellence), then the incentive
compatibility of the Vickrey auction suggests that institutions should
encourage participants to reveal their full depth---to ``bid'' their
best work without strategic withholding.  Academic peer review, at its
best, approximates this ideal: researchers submit their deepest work for
evaluation by experts.  But peer review also suffers from well-known
failures (bias, conservatism, the Matthew effect) that distort the
depth-based selection process.  The mechanism design framework provides
tools for diagnosing and correcting these failures.\footnote{For mechanism
design applied to scientific institutions, see Partha Dasgupta and
Paul David, ``Toward a New Economics of Science,'' \textit{Research
Policy} 23:5 (1994), pp.\ 487--521.}
\end{remark}

\begin{theorem}[Void Composition Preserves Bid, Right (14.8)]
\label{thm:void-bid-right}
$\operatorname{bidValue}(a \compose \void) = \operatorname{bidValue}(a)$.
\end{theorem}

\begin{proof}[Proof sketch]
$a \compose \void = a$ by the void axiom.
\end{proof}

\begin{theorem}[Void Composition Preserves Bid, Left (14.9)]
\label{thm:void-bid-left}
$\operatorname{bidValue}(\void \compose a) = \operatorname{bidValue}(a)$.
\end{theorem}

\begin{proof}[Proof sketch]
$\void \compose a = a$ by the void axiom.
\end{proof}

Adding ``nothing'' to a gift does not increase its prestige.  Padding a
potlatch offering with empty gestures is structurally transparent---the
audience sees through it.  Prefacing a gift with silence (the dramatic pause
before the big reveal) adds nothing to the gift's structural depth.
Ceremony may have psychological effects, but it is orthogonal to substance
in the depth metric.

\begin{example}[Medieval Relic Trade as Prestige Auction]
\label{ex:relics}
The medieval European trade in saints' relics provides a vivid historical
example of the ideatic auction.  From the early Middle Ages through the
Reformation, relics---bones, clothing, instruments of martyrdom, even
phials of saints' blood---were objects of intense competition among
churches, monasteries, and secular rulers.  A relic's ``bid value'' was
its depth: the proximity of the relic to the original saint, the
reliability of its provenance chain (\textit{translatio}), the number
and quality of miracles attributed to it, and the liturgical traditions
that had accumulated around it.

The most prestigious relics---the Crown of Thorns acquired by Louis~IX
for the Sainte-Chapelle, the relics of St.~James at Compostela, the
Holy Blood of Bruges---were ``deep'' in precisely our technical sense:
they were embedded in rich networks of narrative, ritual, and institutional
memory.  Theorem~\ref{thm:void-bid} explains why a church without relics
had zero prestige in the medieval devotional economy: it had bid void.
Theorem~\ref{thm:composed-bid} explains why combining multiple minor
relics (a practice known as \textit{relic bundling}) produced prestige at
most equal to the sum of the individual relics' values---never more.

The relic trade also illustrates Theorem~\ref{thm:winner-bound}: the
prestige impact of even the most extraordinary relic was bounded by the
receiving community's capacity to appreciate it.  A relic of the True
Cross placed in a remote village chapel produced less cultural impact
than the same relic placed in a great cathedral with an educated clergy
capable of composing elaborate liturgical celebrations around it.  The
relic's depth was constant, but $\depth(r)$---the receiver's depth---varied
enormously.\footnote{On the relic trade, see Patrick Geary, \textit{Furta
Sacra} (1978); and Charles Freeman, \textit{Holy Bones, Holy Dust: How
Relics Shaped the History of Medieval Europe} (2011).}
\end{example}

% ----------------------------------------------------------------
\section{Interpretation of Winning Bids}
\label{sec:ch14-interp}
% ----------------------------------------------------------------

\begin{theorem}[Winner's Interpretation Depth Bound (14.10)]
\label{thm:winner-bound}
The depth of the auction interpretation is bounded:
\[
  \depth(\operatorname{auctionInterpret}(\text{bids}, \text{fallback}, r))
  \;\le\;
  \depth(r) + \depth(\operatorname{maxDepthIdea}(\text{fallback}, \text{bids})).
\]
\end{theorem}

\begin{proof}[Proof sketch]
By subadditivity applied to $r \compose \operatorname{maxDepthIdea}(\cdots)$.
\end{proof}

The cultural impact of the winning bid is bounded by the receiver's
background plus the bid's depth.  Even the most extravagant potlatch is
limited by the audience's capacity to appreciate it.  Bourdieu's insight
that cultural capital constrains reception is formalized: a peasant
watching a court ballet and a courtier watching the same performance
produce interpretations of vastly different depth, not because the ballet
differs but because $\depth(r)$ differs.\footnote{Bourdieu's analysis of
aesthetic judgment in \textit{La Distinction} shows that reception of
cultural goods is stratified by the receiver's accumulated cultural
capital---our $\depth(r)$.}

\begin{example}[Academic Citation as Depth-Based Bidding]
\label{ex:citation}
The academic publication system provides a modern instance of depth-based
auction dynamics.  Each published paper is a ``bid'' in the marketplace
of ideas, with its depth determined by the sophistication of its
arguments, the breadth of its evidence, and the novelty of its
contributions.  The ``auction'' takes place through peer review (the
selection mechanism) and citation (the post-auction valuation).

A paper's bid value ($\operatorname{bidValue}(a) = \depth(a)$) reflects
its intellectual complexity.  Papers with greater depth tend to win the
``auction'' of attention: they are published in more prestigious venues,
cited more frequently, and exert greater influence on subsequent research.
Theorem~\ref{thm:void-bid} captures a basic truth of academic life: a
researcher who publishes nothing has zero academic prestige.
Theorem~\ref{thm:composed-bid} applies to collaborative papers: the
depth of a co-authored paper is at most the sum of the co-authors'
individual depths, though it may be less due to coordination costs and
redundancy.

Theorem~\ref{thm:winner-bound} is particularly illuminating: the impact
of a paper on a reader is bounded by the reader's own depth.  A
groundbreaking paper in algebraic geometry will have enormous impact on a
reader with deep mathematical background ($\depth(r)$ large) but minimal
impact on a reader without the prerequisites ($\depth(r)$ small).  This
is the formal basis of the ``prerequisites problem'' in education: students
cannot absorb material that exceeds their background depth, regardless of
the material's intrinsic quality.\footnote{On the sociology of academic
citation, see Robert K.\ Merton, ``The Matthew Effect in Science,''
\textit{Science} 159:3810 (1968), pp.\ 56--63; and Pierre Bourdieu,
\textit{Homo Academicus} (1984).}
\end{example}

\begin{theorem}[Empty Auction Interpretation (14.11)]
\label{thm:empty-interp}
With no bids, the interpretation is the receiver composed with the fallback:
$\operatorname{auctionInterpret}([], \text{fallback}, r) = r \compose \text{fallback}$.
\end{theorem}

\begin{proof}[Proof sketch]
$\operatorname{maxDepthIdea}(\text{fallback}, []) = \text{fallback}$.
\end{proof}

When no external cultural force competes, the receiver interprets only
through their default framework.  Cultural isolation means self-referential
interpretation: the audience, receiving no new bids, composes its background
with tradition alone.

\begin{theorem}[Void Receiver Gets Raw Winner (14.12)]
\label{thm:void-receiver}
If the receiver's background is void, the interpretation is just the winning
bid:
$\operatorname{auctionInterpret}(\text{bids}, \text{fallback}, \void) =
\operatorname{maxDepthIdea}(\text{fallback}, \text{bids})$.
\end{theorem}

\begin{proof}[Proof sketch]
$\void \compose w = w$ by the void axiom.
\end{proof}

The ``blank slate'' receiver---the newborn, the cultural outsider---simply
absorbs the winning idea without transformation.  Cultural naivety means
unmediated reception.  This is the anthropological equivalent of ``going
native'': the observer without pre-existing cultural categories receives the
dominant cultural signal raw, unfiltered by prior interpretive
frameworks.\footnote{The trope of the ethnographer ``going native'' (from
Malinowski onward) captures precisely this: the field worker who sheds their
own cultural background ($r \to \void$) receives the host culture's signals
without the mediation of home categories.}

\begin{example}[Social Media ``Likes'' as Micro-Auctions]
\label{ex:social-media}
Social media platforms implement a continuous, decentralized version of
the ideatic auction.  Each post is a bid for attention, and the
platform's algorithmic feed serves as the auctioneer, selecting the
``deepest'' (most engaging) content for display.  The ``likes,''
``shares,'' and ``retweets'' that follow are a form of post-auction
valuation: the audience registers its interpretation of the winning bids.

The depth-as-prestige identification
(Definition~\ref{def:bid-value}) maps directly onto the platform's
engagement metrics: content with greater depth (complexity, novelty,
emotional resonance) tends to attract more engagement, which the
platform interprets as higher value.  Theorem~\ref{thm:void-bid}
explains why users who post nothing---``lurkers''---accumulate zero social
media prestige.  Theorem~\ref{thm:void-loses} explains why any post, no
matter how trivial, outperforms silence in the attention economy.

However, the social media auction also reveals tensions in our framework.
Platform algorithms optimize for \emph{engagement} rather than
\emph{depth}, and the two do not always coincide: shallow, emotionally
provocative content may generate more engagement than deep, nuanced
analysis.  This suggests that real-world auctions may use metrics other
than depth for bid evaluation---a limitation of the pure depth-as-prestige
identification that our formalization makes explicit.\footnote{For the
attention economy, see Tim Wu, \textit{The Attention Merchants} (2016);
and Zeynep Tufekci, ``YouTube, the Great Radicalizer,'' \textit{The New
York Times} (March 10, 2018).}
\end{example}

\begin{remark}[The Limits of Depth-as-Prestige Identification]
\label{rem:depth-prestige}
Definition~\ref{def:bid-value} identifies bid value with depth:
$\operatorname{bidValue}(a) = \depth(a)$.  This identification is the
central modeling choice of this chapter, and it deserves scrutiny.  In
many cultural contexts, the identification is apt: prestige accrues to
complexity, sophistication, and elaboration.  The Kwakiutl potlatch,
the medieval cathedral, the academic monograph---all derive their
prestige from depth.

But there are important counterexamples.  Zen Buddhism values simplicity
over complexity; the haiku achieves prestige through compression rather
than elaboration; minimalist art (Judd, Andre, Flavin) stakes its claim
on the refusal of depth.  In these contexts, prestige may be inversely
related to depth, or orthogonal to it.  Our formalization captures only
the ``maximalist'' regime in which deeper is more prestigious.  A more
general framework might parameterize the bid-value function, allowing
$\operatorname{bidValue}(a) = f(\depth(a))$ for some monotone (or even
non-monotone) function~$f$, thereby accommodating both the potlatch
(where $f$ is increasing) and the haiku (where $f$ might be decreasing
above a threshold).
\end{remark}

% ----------------------------------------------------------------
\section{Structural Auction Theorems}
\label{sec:ch14-structural}
% ----------------------------------------------------------------

\begin{theorem}[Iterated Bid Depth Bound (14.13)]
\label{thm:iterated-bid}
Repeating a bid $n$ times via composition yields depth at most
$n \cdot \operatorname{bidValue}(a)$:
\[
  \operatorname{bidValue}(\compN{a}{n}) \;\le\; n \cdot \operatorname{bidValue}(a).
\]
\end{theorem}

\begin{proof}[Proof sketch]
By the iterated composition depth bound.
\end{proof}

Repetitive gift-giving---annual potlatch cycles, recurring tribute---has
bounded cumulative prestige.  You cannot achieve infinite status through mere
repetition.  The subadditivity of depth imposes diminishing returns:
the $n$-th iteration of a potlatch adds at most $\depth(a)$, and the
total is bounded by $n \cdot \depth(a)$.

\begin{theorem}[Atomic Bid Bound (14.14)]
\label{thm:atomic-bid}
Atomic ideas have bid value at most one: if $a$ is atomic, then
$\operatorname{bidValue}(a) \le 1$.
\end{theorem}

\begin{proof}[Proof sketch]
By the depth bound on atomic ideas.
\end{proof}

The simplest, most elemental cultural offerings---a single word, a basic
gesture, a minimal ritual act---have minimal prestige value.  Prestige
comes from complexity: the elaborate ceremony outranks the simple nod.
This formalizes the anthropological observation that prestige goods are
invariably \emph{complex} goods: the jade celt, the knotted quipu, the
illuminated manuscript---all derive their prestige from depth, not from
mere existence.

\begin{theorem}[Resonant Bids Produce Resonant Interpretations (14.15)]
\label{thm:resonant-bids}
If two bids resonate, $b_1 \res b_2$, any common receiver produces resonant
interpretations:
\[
  b_1 \res b_2
  \;\Longrightarrow\;
  (r \compose b_1) \res (r \compose b_2).
\]
\end{theorem}

\begin{proof}[Proof sketch]
By left-composition preservation of resonance.
\end{proof}

Competing ritual systems that resonate---e.g., Catholic and Orthodox
Christianity---produce resonant experiences in the same audience.  The laity
perceives structural similarity even across competing institutions.  This
is the formal basis of religious tolerance: if the competing bids resonate,
the audience's interpretations will also resonate, reducing the perceived
distance between rival systems.

\begin{theorem}[Resonant Receivers Agree on Winner (14.16)]
\label{thm:resonant-receivers}
If two receivers resonate, $r_1 \res r_2$, and evaluate the same bid, their
interpretations resonate:
\[
  r_1 \res r_2
  \;\Longrightarrow\;
  (r_1 \compose b) \res (r_2 \compose b).
\]
\end{theorem}

\begin{proof}[Proof sketch]
By right-composition preservation of resonance.
\end{proof}

Members of the same culture watching the same potlatch produce resonant
evaluations.  Cultural consensus on prestige is a structural consequence of
shared background.  This is why prestige hierarchies are culture-specific:
a Kwakiutl potlatch and a Parisian art auction assign prestige to different
bids, but within each culture, members with resonant backgrounds agree on
what is prestigious.

\begin{theorem}[Bid Associativity (14.17)]
\label{thm:bid-assoc}
Composing bids is associative:
$(a \compose b) \compose c = a \compose (b \compose c)$.
\end{theorem}

\begin{proof}[Proof sketch]
By associativity of the monoid.
\end{proof}

In a multi-stage potlatch (gifts given over successive days), the grouping of
days does not affect the total composition.  The ``grand total'' of a
festival is invariant under scheduling---structure trumps sequence.  Whether
the chief gives blankets on Monday and copper shields on Tuesday, or reverses
the order, the total compositional depth is the same (though the individual
days' compositions may differ).

\begin{remark}[Dynamic Auctions and Repeated Games]
\label{rem:dynamic-auctions}
The theorems in this chapter analyze \emph{static} auctions---single-round
competitions in which bids are evaluated once.  But many cultural
competitions are \emph{dynamic}: they unfold over time, with bids placed
in successive rounds and outcomes depending on the history of previous
bids.  The annual potlatch cycle, the multi-year academic tenure process,
the lifelong artistic career---all are repeated auctions in which current
bids are evaluated against the background of past performance.

In a repeated auction, Theorem~\ref{thm:iterated-bid} provides the
fundamental constraint: repeated bidding yields depth at most
$n \cdot \depth(a)$, where $n$ is the number of rounds and $a$ is the
bid.  This linear bound on cumulative prestige has important implications.
It means that an artist who produces works of constant depth~$d$ in each
of $n$ periods accumulates prestige at most $nd$---a linear growth rate
that precludes the exponential prestige accumulation sometimes imagined
in popular accounts of artistic genius.

The theory of repeated games (Aumann, 1959; Friedman, 1971) shows that
repeated interaction can sustain cooperative outcomes not possible in
one-shot games.  In the ideatic setting, this suggests that repeated
cultural competition can sustain \emph{norms} of depth: participants
bid deeply not merely to win the current round but to maintain a reputation
for depth that benefits them in future rounds.  The potlatch chief who
underbids in one year risks being excluded from future exchanges---a
dynamic incentive not captured by our static analysis but consistent with
its structural constraints.\footnote{For repeated game theory, see Drew
Fudenberg and Jean Tirole, \textit{Game Theory} (1991), ch.~5; and for
applications to cultural institutions, see Avner Greif, \textit{Institutions
and the Path to the Modern Economy} (2006).}
\end{remark}

\begin{theorem}[Depth-Zero Bid Composure (14.18)]
\label{thm:zero-depth-comp}
If both bidders have zero depth, their composition has zero depth:
\[
  \depth(a) = 0 \;\wedge\; \depth(b) = 0
  \;\Longrightarrow\;
  \depth(a \compose b) = 0.
\]
\end{theorem}

\begin{proof}[Proof sketch]
By subadditivity, $\depth(a \compose b) \le \depth(a) + \depth(b) = 0$,
and depth is non-negative, so $\depth(a \compose b) = 0$.
\end{proof}

Two culturally empty offerings combined produce nothing.  ``Nothing plus
nothing is nothing''---the formal futility of combining shallow gestures.
Two empty speeches, two blank canvases, two silent prayers: their
composition is as void as each alone.  The potlatch chief who bundles
two worthless gifts has still bid zero.

\begin{lstlisting}[caption={Depth-zero composition in Lean~4}]
theorem zero_depth_composition {a b : I}
    (ha : depth a = 0) (hb : depth b = 0) :
    depth (compose a b) = 0 := by
  have h := depth_compose_le a b
  omega
\end{lstlisting}

% ----------------------------------------------------------------
\section{From Potlatch to Platform: Auctions Across Scales}
\label{sec:ch14-across-scales}
% ----------------------------------------------------------------

The ideatic auction is a structural pattern that recurs across an
extraordinary range of human institutions, from the smallest face-to-face
community to the largest digital platform.  In this section, we trace the
auction pattern across three scales---local, institutional, and
global---showing how the theorems of this chapter apply at each level.

\subsection{Local Auctions: Face-to-Face Prestige Competition}

At the smallest scale, ideatic auctions occur in face-to-face interactions
where individuals compete for interpretive priority through displays of
cultural depth.  The dinner-party conversation, the academic seminar, the
religious disputation---all are local auctions in which participants bid
their ideas and the ``winner'' (the deepest contribution) sets the terms
of subsequent discussion.

Erving Goffman's interaction ritual theory (1967) provides the
micro-sociological grounding for local auctions.\footnote{Erving Goffman,
\textit{Interaction Ritual} (1967); and Randall Collins, \textit{Interaction
Ritual Chains} (2004), which extends Goffman's framework to explain the
accumulation of ``emotional energy'' through successful interaction
rituals.}  In Goffman's framework, each participant in a social encounter
presents a ``face''---a claimed social value---that is evaluated by others.
The participant whose face is most successfully maintained wins the
interaction's ``pot'' of emotional energy and social approval.  In our
terms, each participant's ``face'' is their bid, and the depth of the
bid determines its success.

Theorem~\ref{thm:single-bidder} captures the dynamics of monopolized
conversation: when one participant dominates (a single bidder), they
win the auction by default, provided their depth exceeds the fallback
(the conventional expectation).  Theorem~\ref{thm:void-bid} explains
why silence (void bidding) is socially costly in most interaction
contexts: the participant who says nothing forfeits the competition and
risks ``losing face'' in Goffman's sense.

\subsection{Institutional Auctions: Organized Cultural Competition}

At the institutional scale, auctions are formalized through rules,
procedures, and evaluation criteria.  The academic peer review system,
the art exhibition jury, the literary prize committee, the religious
ordination process---all are institutional auctions in which bids
(papers, artworks, manuscripts, candidates) are evaluated by designated
judges according to established criteria.

Institutional auctions differ from local auctions in several respects
that our framework captures.  First, the evaluation is \emph{delegated}:
the ``receiver'' in Theorem~\ref{thm:winner-bound} is not the general
public but a specialized panel with (typically) high $\depth(r)$.  This
delegation increases the interpretive depth of the evaluation but also
introduces potential biases (the panel's particular composition shapes
which bids are recognized as ``deep'').

Second, institutional auctions are often \emph{iterated}: the Nobel
Prize is awarded annually, tenure reviews occur on a fixed schedule,
art exhibitions recur seasonally.  Theorem~\ref{thm:iterated-bid}
constrains the cumulative prestige of iterated bidding, while the
repeated-game dynamics discussed in Remark~\ref{rem:dynamic-auctions}
explain why institutional participants invest in reputation-building
over time.

Third, institutional auctions impose \emph{format constraints}: papers
must follow certain conventions, artworks must fit exhibition spaces,
candidates must meet eligibility criteria.  These constraints function
as filters on the space of permissible bids, potentially excluding
high-depth contributions that do not conform to institutional norms.
The tension between format constraints and depth maximization is a
perennial source of institutional controversy---from the Salon des
Refusés (1863) to the replication crisis in psychology.

\subsection{Global Auctions: Platforms and Markets}

At the global scale, digital platforms have created auction mechanisms
of unprecedented scope and speed.  Google's search algorithm ranks
web pages by a depth-like metric (PageRank, which measures the
structural centrality of a page in the link graph).  Amazon's
recommendation system selects products through a form of collaborative
filtering that implicitly auctions items for user attention.  TikTok's
algorithmic feed is a continuous auction in which short video ``bids''
compete for screen time based on engagement metrics that proxy for depth.

These global auctions share the structural properties formalized in our
theorems.  Void content receives zero engagement
(Theorem~\ref{thm:void-bid}).  Deeper content tends to outperform
shallower content (Theorem~\ref{thm:deeper-void}).  The impact of
content on any given user is bounded by that user's background
(Theorem~\ref{thm:winner-bound}).  And the composition of content
(playlists, reading lists, curated feeds) is subadditive
(Theorem~\ref{thm:composed-bid}).

The global auction scale also introduces phenomena that push against our
framework's assumptions.  Network effects can create ``viral'' content
whose apparent depth is inflated by social contagion rather than
intrinsic complexity.  Algorithmic curation may optimize for engagement
rather than depth, creating incentives for ``clickbait''---content
designed to exploit cognitive biases rather than demonstrate genuine
cultural depth.  These phenomena suggest extensions of the framework
in which the bid-value function incorporates social and algorithmic
amplification factors beyond the intrinsic depth of the idea itself.

\subsection{Cross-Scale Connections}

The three scales of auction---local, institutional, and global---are not
independent but mutually reinforcing.  An academic's success in local
auctions (seminars, conferences) builds the reputation that enables
success in institutional auctions (publications, grants, prizes), which
in turn amplifies their presence in global auctions (media coverage,
public intellectual status).  This cross-scale amplification is consistent
with our framework: each level of auction is subject to the same
structural constraints, and success at one level provides the depth
(accumulated cultural capital) that enables competitive bidding at the
next level.

The cross-scale perspective also illuminates the historical transition
from potlatch to platform.  The potlatch was a local auction embedded
in a specific community; its prestige effects were geographically
bounded.  The modern art market is an institutional auction with global
reach; its prestige effects span continents.  The attention economy of
social media is a global auction operating in real time; its prestige
effects are instantaneous and worldwide.  At each stage, the auction's
scope has expanded, but the structural constraints---subadditivity,
void loss, bounded reception---remain invariant.  The potlatch chief
and the TikTok influencer are subject to the same ideatic algebra.

% ----------------------------------------------------------------
\section{Conclusion}
\label{sec:ch14-conclusion}
% ----------------------------------------------------------------

The ideatic auction formalizes a pervasive pattern in human cultural
life: competition for interpretive priority through depth.  The
structural results of this chapter---void always loses
(Theorem~\ref{thm:void-bid}), composition is subadditive
(Theorem~\ref{thm:composed-bid}), reception is bounded by audience
capacity (Theorem~\ref{thm:winner-bound}), and zero-depth bids compose to
zero (Theorem~\ref{thm:zero-depth-comp})---impose hard constraints on
the ritual economy.

These constraints connect to three major anthropological traditions.
Mauss's gift theory is formalized through the obligation to bid non-void
(Theorem~\ref{thm:void-bid}); Veblen's conspicuous consumption is
captured by the depth-as-prestige identification
(Definition~\ref{def:bid-value}); and Bourdieu's symbolic capital is
reflected in the receiver-bounded interpretation
(Theorem~\ref{thm:winner-bound}).

In the final chapter of this sequence, we turn from competitive bidding
to \emph{curation}: the receiver's active management of their own
repertoire, connecting to museum theory, education policy, and the
political economy of knowledge.
